\chapter{Generative Adversarial Networks --- GANs From Scratch}
\label{chap:ch13}

\section{The Intuition of Adversarial Competition}
Generative Adversarial Networks (Goodfellow et al., Isola et al.~\cite{isola2017pix2pix}) formulate deep learning as a two-player zero-sum game between two competing neural networks:
\begin{enumerate}
    \item \textbf{The Generator ($\mathcal{G}$):} Plays the role of an art forger attempting to paint realistic clear-sky optical satellite imagery beneath clouds.
    \item \textbf{The Discriminator ($\mathcal{D}$):} Plays the role of an art detective examining satellite images to determine whether a given patch is a genuine satellite acquisition or an artificial synthetic reconstruction.
\end{enumerate}

\begin{conceptbox}{Why Adversarial Competition Eliminates Blurring}
If a standard generator produces a blurry patch to safely minimize $L_1$ error, a conventional loss function will not complain. But a trained Discriminator will instantly recognize: \emph{"Real satellite sensors never produce blurry patches with zero high-frequency noise; this is obviously fake!"} and assign it a punitive error score. To fool the detective, the generator is forced to synthesize authentic high-frequency textures, crop field rows, and natural edge gradients!
\end{conceptbox}

\section{Conditional PatchGAN Architecture}
In standard GANs, the discriminator inspects the entire image and outputs a single global probability scalar between 0 and 1.

In \textbf{PatchGAN} (Isola et al.~\cite{isola2017pix2pix}), the discriminator is fully convolutional: its final layer produces a spatial grid of logit predictions (in our model, a $16 \times 16$ grid). Each neuron in this grid evaluates the realism of an overlapping local spatial patch ($\approx 70\text{ m} \times 70\text{ m}$ receptive field). This forces the discriminator to focus strictly on local high-frequency textures rather than global layout.

\section{The Verified 11-Channel High-Pass Conditioned Discriminator}
In \texttt{losses/patchgan\_loss.py} (lines 15--75), our project introduces a novel discriminator conditioning scheme specifically engineered to combat over-smoothing.

Figure~\ref{fig:patchgan_11ch} illustrates the exact architecture and tensor flow.

\begin{figure}[htbp]
\centering
\begin{tikzpicture}[
    font=\footnotesize\sffamily,
    chan/.style={draw=black!75, thick, rounded corners=3pt, minimum width=3.8cm, minimum height=0.6cm, align=center},
    layer/.style={draw=purple!80!black, thick, fill=purple!12, rounded corners=4pt, minimum width=4.4cm, minimum height=0.8cm, align=center},
    arr/.style={draw=black!85, thick, -Stealth}
]

% Inputs
\node[chan, fill=blue!15] (c1) at (0, 4.2) {Base Optical: $3.0 \times \mathbf{Y}$ ($4$ ch)};
\node[chan, fill=red!15] (c2) at (0, 3.4) {High-Pass: $10.0 \times (\mathbf{Y} - \text{AvgPool}_{5\times 5}(\mathbf{Y}))$ ($4$ ch)};
\node[chan, fill=cyan!15] (c3) at (0, 2.6) {Sentinel-1 SAR Condition: $\mathbf{X}_{S1}$ ($2$ ch)};
\node[chan, fill=orange!15] (c4) at (0, 1.8) {Soft Cloud Mask: $\mathbf{M}$ ($1$ ch)};

% Concat node
\node[draw=black, thick, fill=gray!20, rounded corners=4pt, align=center, minimum width=4.2cm, minimum height=0.8cm] (cat) at (0, 0.6) {\textbf{11-Channel Stacked Tensor}\\\scriptsize $[B, 11, 128, 128]$};

\draw[arr] (c1) -- (cat);
\draw[arr] (c2) -- (cat);
\draw[arr] (c3) -- (cat);
\draw[arr] (c4) -- (cat);

% Discriminator layers
\node[layer] (l1) at (6.0, 3.8) {\texttt{SNConv2d(11, 64, k=4, s=2)}\\\scriptsize $[B, 64, 64, 64]$ + LeakyReLU($0.2$)};
\node[layer] (l2) at (6.0, 2.6) {\texttt{SNConv2d(64, 128, k=4, s=2)}\\\scriptsize $[B, 128, 32, 32]$ + LeakyReLU($0.2$)};
\node[layer] (l3) at (6.0, 1.4) {\texttt{SNConv2d(128, 256, k=4, s=2)}\\\scriptsize $[B, 256, 16, 16]$ + LeakyReLU($0.2$)};
\node[layer] (l4) at (6.0, 0.2) {\texttt{SNConv2d(256, 512, k=4, s=1)}\\\scriptsize $[B, 512, 16, 16]$ + LeakyReLU($0.2$)};
\node[layer, fill=green!15] (l5) at (6.0, -1.0) {\texttt{SNConv2d(512, 1, k=4, s=1)}\\\scriptsize Output Logits: $[B, 1, 16, 16]$};

\draw[arr] (cat.east) -| (3.5, 3.8) -- (l1.west);
\draw[arr] (l1) -- (l2);
\draw[arr] (l2) -- (l3);
\draw[arr] (l3) -- (l4);
\draw[arr] (l4) -- (l5);

% Patch Weighting
\node[draw=orange!80!black, thick, fill=orange!12, rounded corners=4pt, align=center, minimum width=4.4cm, minimum height=1.0cm] (wloss) at (6.0, -2.4) {\textbf{Mask-Weighted Hinge Loss}\\$w_{\text{patch}} = 0.25 + 0.75 \times \text{AvgPool}_{8\times 8}(\mathbf{M})$\\$\mathcal{L}_D = \text{mean}(w_{\text{patch}} \cdot \text{ReLU}(1 \pm D(\cdot)))$};

\draw[arr] (l5) -- (wloss);

\end{tikzpicture}
\caption{Architecture of the 11-channel High-Pass Conditioned PatchGAN Discriminator (\texttt{losses/patchgan\_loss.py}, lines 15--75). The discriminator receives optical imagery (scaled by $3.0$), high-pass residual detail (scaled by $10.0$), SAR conditions, and the cloud mask. Every convolution uses Spectral Normalization to prevent mode collapse.}
\label{fig:patchgan_11ch}
\end{figure}


\subsection{Exact Channel Breakdown: 11 Input Channels}
The discriminator receives an 11-channel concatenated tensor:
\begin{enumerate}
    \item \textbf{Base Optical Image ($4$ channels):} $3.0 \times \mathbf{Y}$ (B04, B03, B02, B08 scaled for optimal gradient response).
    \item \textbf{High-Pass Residual Detail ($4$ channels):}
    \begin{equation}
    \mathbf{R}_{\text{HP}} = 10.0 \times \big(\mathbf{Y} - \text{AvgPool}_{5\times 5}(\mathbf{Y})\big)
    \end{equation}
    \item \textbf{Sentinel-1 SAR Condition ($2$ channels):} $\mathbf{X}_{\text{SAR}}$ (VV and VH polarizations).
    \item \textbf{Soft Cloud Mask ($1$ channel):} $\mathbf{M} \in [0, 1]$.
\end{enumerate}
Total channels: $4 + 4 + 2 + 1 = \mathbf{11\text{ channels}}$.

\begin{mathbox}{How the High-Pass Filter ($5 \times 5$ Residual) Works}
Subtracting a local $5 \times 5$ average pool removes all low-frequency macroscopic colors, leaving only pure high-frequency edges, noise, and textures:
\[
\text{Low-Frequency Base} = \text{AvgPool}_{5\times 5}(\mathbf{Y})
\]
\[
\text{High-Frequency Detail} = \mathbf{Y} - \text{AvgPool}_{5\times 5}(\mathbf{Y})
\]
Multiplying this residual by $10.0$ magnifies subtle texture differences by an order of magnitude. If the generator produces an over-smoothed, blurry reconstruction, its high-pass channels are nearly zero, allowing the discriminator to instantly penalize it!
\end{mathbox}

\subsection{Spectral Normalization}
Every convolutional layer in the discriminator is wrapped with \textbf{Spectral Normalization} (\texttt{torch.nn.utils.spectral\_norm}, Miyato et al.~\cite{miyato2018spectral}). Spectral normalization divides weight matrices by their largest singular value (spectral norm $\sigma(\mathbf{W})$), strictly enforcing a 1-Lipschitz continuity constraint:
\begin{equation}
\|\mathcal{D}(\mathbf{x}_1) - \mathcal{D}(\mathbf{x}_2)\| \le 1.0 \cdot \|\mathbf{x}_1 - \mathbf{x}_2\|
\end{equation}
This mathematically guarantees that the discriminator's gradients cannot explode or vanish, stabilizing adversarial training and completely eliminating mode collapse!

\subsection{Mask-Weighted Hinge Adversarial Loss}
The discriminator loss is weighted by the downsampled cloud mask so that adversarial training concentrates on reconstructed cloud regions while applying baseline supervision to clear regions:
\begin{equation}
w_{\text{patch}} = 0.25 + 0.75 \times \text{AvgPool}_{8\times 8}(\mathbf{M})
\end{equation}
The Hinge loss formulation is computed as:
\begin{align}
\mathcal{L}_D &= \text{mean}\Big( w_{\text{patch}} \odot \big[ \text{ReLU}(1.0 - \mathcal{D}(\mathbf{X}_{\text{real}})) + \text{ReLU}(1.0 + \mathcal{D}(\mathbf{X}_{\text{fake}})) \big] \Big) \\
\mathcal{L}_{G,\text{adv}} &= \text{mean}\Big( w_{\text{patch}} \odot \big[ -\mathcal{D}(\mathbf{X}_{\text{fake}}) \big] \Big)
\end{align}

\begin{researchbox}{Active Status of the GAN in the Repository}
We audited all training scripts to establish the exact factual status of the GAN:
\begin{itemize}
    \item \textbf{Stage 1 (\texttt{train\_ecrformer.py}):} GAN was \textbf{INACTIVE}. Trained with baseline SDFL loss only. Reconstructions suffered from over-smoothing (detail ratio $23\%\text{--}47\%$).
    \item \textbf{Stage 4 (\texttt{train\_ecrformer\_v4\_gan.py}):} The 11-channel PatchGAN was \textbf{FULLY ACTIVE} and trained in an alternating optimization loop. This run directly generated the project's premier checkpoint: \texttt{ecrformer\_sweetspot\_best/best\_model.pth} (retaining $85.76\%$ detail ratio and $29.97\text{ dB}$ PSNR on heavy clouds)!
\end{itemize}
\end{researchbox}

\section{Summary of Chapter 13}
\begin{enumerate}
    \item GANs create an adversarial game between a generator and discriminator that forces generated imagery to exhibit natural high-frequency textures.
    \item PatchGAN evaluates local $16 \times 16$ spatial grids, providing localized gradient feedback.
    \item The 11-channel discriminator combines base optical, $10\times$ magnified high-pass residual, SAR conditions, and cloud masks.
    \item Spectral normalization enforces 1-Lipschitz continuity, eliminating mode collapse.
    \item The GAN was fully active during Stage 4 training, producing the verified sweetspot checkpoint.
\end{enumerate}

\section{Practice Exercises}
\begin{enumerate}
    \item Calculate the total number of channels entering the discriminator when base optical has 4 channels, high-pass has 4 channels, SAR has 2 channels, and mask has 1 channel.
    \item Why does Spectral Normalization prevent discriminator gradients from exploding?
    \item What is the value of $w_{\text{patch}}$ over a completely clear region ($\mathbf{M} = 0.0$)? Over a completely cloudy region ($\mathbf{M} = 1.0$)?
\end{enumerate}

\subsection*{Solutions}
\begin{enumerate}
    \item $4 + 4 + 2 + 1 = 11$ channels.
    \item Spectral normalization bounds the Lipschitz constant of the network by scaling each weight matrix by its spectral norm ($\mathbf{W} / \sigma(\mathbf{W})$). This guarantees that the magnitude of the gradient $\|\nabla_{\mathbf{x}} \mathcal{D}(\mathbf{x})\| \le 1.0$ everywhere, preventing the discriminator from overpowering the generator with explosive gradients.
    \item Over clear regions: $w_{\text{patch}} = 0.25 + 0.75(0.0) = 0.25$. Over cloudy regions: $w_{\text{patch}} = 0.25 + 0.75(1.0) = 1.00$. Cloudy patches receive $4\times$ stronger adversarial gradient feedback than clear patches!
\end{enumerate}
